\section{Introduction}\label{sec:intro}

The combinator $S$, with rule $S\,x\,y\,z\to x\,z\,(y\,z)$, never erases a subterm, so on
its own it cannot express $K$~\citep{waldmann-thesis}; whether it is computationally
universal is open. The closed terms built from $S$ alone behave in intricate ways. Some normalize, some
reduce forever, and whether a term reaches a head normal form ($S$, $S\,X$ or $S\,X\,Y$)
depends on the whole term. Waldmann showed that normalization of these terms is
decidable and that the normalizing terms form a regular tree language
\citep{waldmann-thesis,waldmann-ic}; see \citet[Theorem~9]{dots}.
%% Checked 2026-11-20 against the thesis abstract (thesis.pdf, from thesis.dvi) and Dance Thm 9.
\citet{dots} surveyed what is known and posed open questions. Their Question~10 offers a
35-state bottom-up tree automaton (their Fig.~5.14) and asks whether it accepts exactly
the closed $S$-terms with a head normal form. They expect one direction to be easy and
call the other, that every accepted term has a head normal form, the challenge.

We answer the question positively.

\begin{theorem}\label{thm:main}
A closed $S$-term has a head normal form if and only if the automaton of
\textup{\citet[Fig.~5.14]{dots}} accepts it.
\end{theorem}

The proof has two parts. First (Sections~\ref{sec:prelims} and~\ref{sec:q10}), the
automaton is a homomorphic image of a 38-element model of the $S$-law. Its state is
therefore invariant under reduction, the rejecting state is the image of a single
absorbing element, and rejection implies the absence of a head normal form. Second
(Section~\ref{sec:soundness}), head reduction of a term whose model value is not that
element terminates. The termination measure reads a weight along the spine of the term.
A head step changes the spine only in a window of three items, and the model guarantees
that the walk resumes after the window in the same state, so the weight drops whenever a
finite inequality holds on every window. A certificate of small integers found by integer
programming satisfies these inequalities, and two tie-breakers handle the windows where
the weight stays the same.

Every step is mechanized in Lean~4~\citep{lean4} within \texttt{lean-ski}~\citep{leanski}.
The finite facts about the tables are checked by \texttt{decide} in the kernel, without
compiled evaluation, and the axiom gate admits only \texttt{propext} and
\texttt{Quot.sound}. We use the notation of \citet{barendregt}.
